\documentclass[12pt]{article}
\usepackage[margin=1in]{geometry}
\usepackage{amsmath,amssymb,booktabs,tabularx,array,setspace}
\usepackage[authoryear,round]{natbib}
\usepackage[colorlinks=true,linkcolor=blue,citecolor=blue,urlcolor=blue]{hyperref}
\setstretch{1.1}
\setlength{\parindent}{2em}
\renewcommand{\tablename}{Table}

\title{Directional Congestion and Tidal-Lane Reallocation}
\author{
Yige Duan\thanks{Antai College of Economics and Management, Shanghai Jiao Tong University; \href{mailto:yige.duan@sjtu.edu.cn}{yige.duan@sjtu.edu.cn}.}
\and
Xiaoruo Feng\thanks{Peking University HSBC Business School; \href{mailto:xiaoruo.feng@stu.pku.edu.cn}{xiaoruo.feng@stu.pku.edu.cn}.}
\and
Yizhen Gu\thanks{Peking University HSBC Business School; \href{mailto:yizhengu@phbs.pku.edu.cn}{yizhengu@phbs.pku.edu.cn}.}
}
\date{July 2026}

\begin{document}
\maketitle

\begin{abstract}
Urban congestion is directional when peak traffic loads the two directions of the same corridor unequally. We combine national route queries, high-frequency road-segment speeds, and commuting flows to measure this asymmetry, then incorporate directed congestion into a quantitative spatial-equilibrium model. In the top 100 congested Chinese cities, the flow-weighted median directional ratio is 0.878. In Beijing, a tidal-lane counterfactual reallocates one lane from the faster to the slower morning direction on the 100 reciprocal corridors with the greatest flow-weighted excess-delay mass. Holding total corridor capacity fixed, the policy raises the model's short-run welfare index by 0.87 percent and its long-run index by 0.32 percent. Long-run congestion feedback and residence--workplace sorting absorb most of the average route-time saving. A reciprocal-symmetry diagnostic also changes welfare and spatial allocations before any policy shock. Directional costs therefore matter for both the measurement of urban congestion and the evaluation of capacity-reallocation policies.
\end{abstract}

\section{Introduction}

Peak commuting traffic often has a direction. Jobs cluster differently from residences, so the morning demand placed on one side of a corridor can exceed demand on the reverse side. This feature is central to reversible or tidal lanes, which reallocate capacity across directions without increasing the total lane supply. Existing work studies peak-period congestion and reversible-lane operations \citep{vickrey1969,arnott1993,wolshon2006}, while quantitative spatial models typically represent transport frictions through bilateral or network costs \citep{ahlfeldt2015,allenarkolakis2022}. The relevant spatial object is nevertheless directed: if traffic differs between $(k,l)$ and $(l,k)$, then the associated commuting costs and market-access terms differ as well.

This paper develops a measurement and modelling framework for this problem. First, it measures directional congestion at national trip, within-city commuter, and road-segment levels. Second, it embeds directed link capacity, route traffic, and endogenous congestion in a quantitative spatial equilibrium. Third, it evaluates a tidal-lane reallocation in Beijing and compares that policy with capacity expansion. The counterfactuals hold baseline fundamentals fixed and resolve the traffic and location-choice equilibrium. They are therefore model-based policy calculations rather than reduced-form estimates of an operating tidal-lane program.

\section{Data and Directional Congestion}

The empirical inputs combine three objects. National trip-level routes cover township pairs in the top 100 congested Chinese cities in 2025. The within-city speed panel records directed road-segment speeds at ten-minute frequency over May 19--25, 2025. The Beijing implementation further deploys Baidu 2022 home--work commuting flows, road geometry, functional class, and lane records. The model is constructed on a directed adjacent-grid network over $3\times3$ km Beijing cells; each physical road supplies two directed links.

For any reciprocal pair, define the directional ratio as
\begin{equation}
  \mathrm{SAR}_{ij}=\frac{\min\{x_{ij},x_{ji}\}}{\max\{x_{ij},x_{ji}\}},
  \label{eq:sar}
\end{equation}
where $x_{ij}$ is the direction-specific route time or road-segment speed, according to the data source. A ratio of one denotes equality across directions, and lower values denote stronger directional asymmetry. In the national trip-level sample, the flow-weighted median $\mathrm{SAR}$ is 0.878 and 58 percent of pair-flow has $\mathrm{SAR}\leq0.90$ during the morning peak. Higher city-level Baidu congestion is associated with stronger morning asymmetry.

The same pattern appears within cities. In Shenzhen, the median ratio is 0.819 when grid pairs are weighted by commuter exposure; 71.6 percent of 4.4 million commuters travel on pairs with $\mathrm{SAR}\leq0.90$. In Beijing, the median morning ratio across 3,670 two-way centerlines is 0.884, and 54.7 percent of centerlines have $\mathrm{SAR}\leq0.90$. A more stringent screen, $\mathrm{SAR}\leq1/3$, identifies 179 centerlines, 4.9 percent of the sample and about 213 km, as candidate corridors for a directional capacity policy.

Jobs--housing separation supplies an economic source of this directional load. On the Beijing square grid, the Duncan dissimilarity index between jobs and residences is 0.26, relative to a random-allocation floor of 0.004. Morning road commuting is radial inward: 74.5 percent of trips head toward the job-weighted center, and the inbound-to-reverse ratio is 2.9. These facts motivate a model in which directional travel costs reshape both bilateral commuting and the spatial allocation of residents and workplaces.

\section{Directed Spatial Equilibrium}

Let $(\mathcal{N},\mathcal{E})$ denote the directed road network. A physical road contains the two directed links $(k,l)$ and $(l,k)$, with free-flow traversal time $t_{kl}^{ff}$, lane capacity $n_{kl}$, traffic $\Xi_{kl}$, and realised traversal time $t_{kl}$. A worker choosing residence $i$, workplace $j$, and route $r\in\mathcal{R}_{ij}$ receives
\begin{equation}
  V_{ij,r}(\omega)=\frac{u_iw_j}{C_{ij,r}}\varepsilon_{ij,r}(\omega),
  \qquad
  C_{ij,r}=\prod_{\ell=1}^{K_r}t_{r_{\ell-1},r_\ell},
  \label{eq:utility}
\end{equation}
where the idiosyncratic term is independently Fr\'echet distributed with shape parameter $\theta$. Integrating over route choices yields the bilateral commuting-cost index
\begin{equation}
  \tau_{ij}^{-\theta}=\sum_{r\in\mathcal{R}_{ij}}C_{ij,r}^{-\theta}.
  \label{eq:tau}
\end{equation}
Directional congestion generally implies $\tau_{ij}\neq\tau_{ji}$.

With residence and workplace marginals, bilateral flows take the doubly constrained gravity form
\begin{equation}
  L_{ij}=\tau_{ij}^{-\theta}
  \left(\frac{L_i^R}{\Pi_i}\right)
  \left(\frac{L_j^F}{P_j}\right),
  \qquad
  L_i^R=\sum_jL_{ij},\quad L_j^F=\sum_iL_{ij},
  \label{eq:gravity}
\end{equation}
where $\Pi_i$ and $P_j$ are residence-side and workplace-side market-access terms. Residential amenity and workplace productivity depend on local activity:
\begin{equation}
  u_i=\bar u_i(L_i^R)^\beta,
  \qquad
  w_j=\bar A_j(L_j^F)^\alpha.
  \label{eq:fundamentals}
\end{equation}
The aggregate welfare index is the Fr\'echet inclusive-value scale,
\begin{equation}
  \bar W=\left[\sum_{i,j}\tau_{ij}^{-\theta}u_i^\theta w_j^\theta\right]^{1/\theta}.
  \label{eq:welfare}
\end{equation}

Traffic and traversal costs are jointly determined. The implementation anchors counterfactual travel times to observed directed peak speeds:
\begin{equation}
  t_{kl}=\max\!\left\{t_{kl}^{ff},\;
  t_{kl}^{obs}\left[
  \left(\frac{\Xi_{kl}}{\Xi_{kl}^{obs}}\right)
  \left(\frac{n_{kl}^{obs}}{n_{kl}}\right)
  \right]^\lambda\right\}.
  \label{eq:anchor}
\end{equation}
Baseline amenities and productivities are inverted to make observed residence and workplace shares an equilibrium at observed directed travel times. Counterfactuals then hold $(\bar u_i,\bar A_j)$ fixed and jointly resolve edge times, route traffic, bilateral costs, commuting flows, location shares, market access, and welfare.

The main calibration sets $\theta=6.83$, $\alpha=0.05$, $\beta=-0.10$, and $\lambda=1.1$. The first parameter governs the response to commuting costs, $\alpha$ captures workplace agglomeration, $\beta$ captures residential crowding, and $\lambda$ maps traffic per lane into directed traversal time. The cost elasticity follows the spatial-equilibrium calibration in \citet{ahlfeldt2015} and \citet{allenarkolakis2022}; the congestion elasticity follows \citet{barwick2024}.

\section{What a Symmetric Cost Surface Misses}

The initial diagnostic compares the directed baseline with a reciprocal cost surface. It holds the network, lane counts, inverted fundamentals, parameters, solver, and welfare formula fixed, then replaces each reciprocal pair of observed and free-flow times with its geometric mean. This comparison isolates the consequence of imposing reciprocal symmetry on the cost surface.

The symmetric specification raises $\bar W$ by 1.10 percent relative to the directed baseline, changes the mean origin--destination route cost by 0.10 percent, and produces residence and workplace dissimilarities of 3.36 and 6.15 percent, respectively. Under a hub-productivity shock, the symmetric anchored model also attenuates the route-cost response from 1.93 to 1.81 percent. Symmetry thus changes the baseline spatial allocation before a transport-policy shock is introduced.

\section{Tidal-Lane Counterfactuals}

For each two-way corridor, let $s$ denote the slower morning direction and $f$ the faster direction. A tidal-lane policy transfers one lane to the slower direction:
\begin{equation}
  n_s'=n_s+1,
  \qquad
  n_f'=n_f-1,
  \qquad
  n_s+n_f=n_s'+n_f'.
  \label{eq:tidal}
\end{equation}
The primary policy ranks reciprocal corridors by the slow-minus-fast congestion-delay mass,
\begin{equation}
  \mathrm{TCS}=
  \Xi_s(t_s^{obs}-t_s^{ff})-
  \Xi_f(t_f^{obs}-t_f^{ff}),
  \label{eq:tcs}
\end{equation}
and treats the top 100 pairs. A positive score assigns a lane toward the direction with greater baseline flow times excess travel time. The policy leaves total corridor capacity and net lane-kilometres unchanged.

Table~\ref{tab:tidal} reports the main equilibrium results. In the short run, residence and workplace shares remain at their baseline values. The policy reduces the route cost by 1.63 seconds per baseline road commuter, or 2,400 person-minutes per day, and raises $\bar W$ by 0.87 percent. The treated corridors are economically exposed: 56.4 percent of baseline origin--destination commuter mass has shortest paths that pass through them. The slow-direction links save 20.3 percent in the short run.

\begin{table}[htbp]
\centering
\caption{Top-100 tidal-lane reallocation: main results}
\label{tab:tidal}
\begin{tabular}{lcc}
\toprule
Outcome & Short run & Long run \\
\midrule
Welfare index, $\bar W$ & $+0.87\%$ & $+0.32\%$ \\
Route cost, seconds per commuter & $-1.63$ & $+0.06$ \\
Route cost, person-minutes per day & $-2{,}400$ & $+95$ \\
Slow-direction link time change & $-20.3\%$ & $-13.6\%$ \\
\bottomrule
\end{tabular}
\begin{flushleft}\footnotesize
\emph{Notes:} The long-run route-cost change decomposes into the short-run direct effect ($-1.63$ seconds), a congestion-feedback wedge ($+0.99$ seconds), and a sorting wedge ($+0.71$ seconds). Negative route-cost entries denote time savings.
\end{flushleft}
\end{table}

Allowing residence--workplace sorting changes the interpretation of the short-run time gain. Network congestion feedback adds 0.99 seconds per commuter and sorting adds 0.71 seconds, leaving a long-run route-cost change of $+0.06$ seconds per commuter. The welfare index nevertheless remains 0.32 percent above its baseline level because workers value access to the full set of residence--workplace opportunities. Residential shares increase near treated endpoints, whereas employment gains concentrate in fewer locations. Mean log residential shares increase by 0.017 at treated endpoints and fall by 0.012 elsewhere; the corresponding workplace changes are $-0.006$ and $-0.016$.

The policy comparison separates reallocating capacity from adding capacity. A targeted road-expansion policy adds one lane in each direction on the 100 corridors with the highest baseline flow times excess travel time. It raises short-run and long-run welfare by 1.26 and 0.44 percent, respectively, but requires 1,140 additional lane-kilometres. A uniform expansion with a matched 1,092 lane-kilometre budget raises long-run welfare by 0.10 percent. The tidal policy yields its 0.32 percent long-run welfare gain with zero net lane-kilometres.

Targeting criteria also matter. Selecting the top 100 links with a local welfare-elasticity score raises long-run welfare by 0.35 percent, compared with 0.32 percent under the flow-weighted excess-delay score and 0.13 percent under a speed-only rule. The welfare-elasticity score combines the local edge gain with upstream residence access and downstream workplace opportunity. Its rank correlation with the speed-only score is 0.08, whereas the correlation with the flow-times-delay score is 0.65.

\section{Scope and Conclusion}

The main results use the stated calibration and an anchored endogenous-congestion equilibrium. Across the reported $\theta$ grid from 3 to 10, long-run welfare gains remain positive, declining from 0.75 to 0.21 percent. Across the converged $\alpha$--$\beta$ combinations reported in the sensitivity analysis, long-run gains range from 0.28 to 0.37 percent. A retained lower-slope congestion check with $\lambda=0.57$ instead yields a long-run welfare change of $-0.078$ percent. The policy calculation also does not endogenize mode switching, so benefits for travellers who would move from another mode to driving are outside the welfare accounting.

Directional congestion is pervasive in the measurement exercises and materially changes the model's spatial equilibrium. The Beijing counterfactual indicates that a targeted tidal-lane reallocation can raise the welfare index without expanding the road network, although congestion feedback and sorting substantially alter the accompanying route-time effect. Capacity policies should therefore be evaluated on a directed network and with spatial adjustment, rather than through a reciprocal commuting-cost surface or a treated-segment speed change alone.

\bibliographystyle{apalike}
\bibliography{../../PreSlides/references}

\end{document}
